\section{Conclusions and  Future  Directions}\label{sec:conclusion}

Multi-agent RL has long been an active and significant research area in reinforcement learning, in view of the  ubiquity of  sequential decision-making with multiple agents coupled in their actions and information. In stark contrast to its great empirical success, theoretical understanding of MARL algorithms is well recognized  to be challenging and relatively lacking in the literature. Indeed,  establishing an encompassing theory for MARL requires tools spanning  dynamic programming, game theory, optimization theory, and statistics, which are  non-trivial  to unify and investigate within one context. 

In this chapter, we have provided a selective overview of mostly recent MARL algorithms, backed by   theoretical analysis, followed by several high-profile but challenging applications that have been addressed lately. Following  the classical overview \cite{bu2008comprehensive}, we have categorized the algorithms into three groups: those solving problems that are  fully cooperative, fully competitive, and a mix of the two.  Orthogonal to the existing reviews on MARL, this chapter has laid emphasis on several new angles and taxonomies of MARL theory,  some of which have been drawn  from our own research endeavors and interests. We note that our overview  should not be viewed  as a  comprehensive one, but instead as a focused one dictated by our own interests and expertise,  which should appeal to  researchers of similar interests, and provide  a stimulus for  future research directions in this general topical area.  Accordingly, we have identified 
 the following paramount while open  avenues for future research on MARL theory.   
 
\vspace{6pt}
\noindent{\bf Partially observed settings:} 
\vspace{3pt} 
Partial observability of the system states and the actions of other agents is quintessential and inevitable in many practical MARL applications. In general, these settings can be modeled as a partially observed stochastic game (POSG), which includes the cooperative setting with a common reward function, i.e., the Dec-POMDP model, as a special case. Nevertheless, as pointed out in \S\ref{subsec:partially_observed}, even the cooperative task is NEXP-complete  \cite{bernstein2002complexity} and difficult to solve. In fact, the information state for optimal decision-making in POSGs can be  very complicated  and involve  belief generation over the opponents' policies 
\cite{hansen2004dynamic}, compared to that  in POMDPs, which requires belief on only states. This difficulty essentially stems  from the heterogenous beliefs of agents  resulting from their own observations obtained from the model, 
an inherent challenge of MARL  mentioned  in \S\ref{sec:challenges} due to various information structures. It might be possible to start by  generalizing the centralized-learning-decentralized-execution scheme for solving Dec-POMDPs \cite{amato2015scalable,dibangoye2018learning} to solving POSGs. 
%Indeed, POSGs can be equivalently transformed into a virtual \emph{central game} over some information states, under certain conditions \cite{nayyar2013common,}.

\vspace{6pt}
\noindent{\bf Deep MARL theory:} 
\vspace{3pt} 
As mentioned in \S\ref{subsec:scala_issue}, using deep neural networks for function approximation can address the scalability issue in MARL. In fact, most of the recent empirical successes in MARL result from the use of DNNs \cite{heinrich2016deep,lowe2017multi,foerster2017counterfactual,gupta2017cooperative,omidshafiei2017deep}. Nonetheless, because of  lack of theoretical backings, we have not included details of these algorithms in this chapter. Very recently, a few attempts have been made to understand the global convergence of several single-agent deep RL algorithms, such as  neural TD learning \cite{cai2019neural} and  neural policy optimization  \cite{wang2019neural,liu2019neural},  when overparameterized neural networks \cite{arora2018optimization,li2018learning} are used. It is thus promising to extend these results to multi-agent  settings, as initial steps toward theoretical understanding of deep MARL. 

\vspace{6pt}
\noindent{\bf Model-based MARL:} 
\vspace{3pt} 
It may be slightly surprising that very few MARL algorithms in the literature are \emph{model-based}, in the sense that   the MARL model is first estimated, and then used as a nominal one to design algorithms. To the best of our knowledge, 
the only existing model-based MARL  algorithms include the early one in \cite{brafman2000near} that solves single-controller-stochastic games, a special zero-sum MG; and the later improved one in \cite{brafman2002r}, named R-MAX, for zero-sum MGs. These  algorithms are also built upon the principle of {optimism in the  face of uncertainty} \cite{auer2007logarithmic, jaksch2010near}, as several aforementioned model-free ones. Considering recent progresses in model-based RL, especially its provable advantages over model-free ones in certain regimes  \cite{tu2018gap,sun2019model}, it is worth generalizing these results to MARL to improve its sample efficiency. 
 
\vspace{6pt}
\noindent{\bf Convergence of policy gradient methods:}
\vspace{3pt}
As mentioned in  \S\ref{subsec:mixed_setting}, the convergence result of vanilla policy gradient method in general MARL is  mostly negative, i.e., it may avoid even the local NE points in many cases.   
This is essentially related to the challenge of non-stationarity in MARL, see \S\ref{subsec:non_stationary}. Even though some remedies have been advocated  \cite{balduzzi2018mechanics,letcher2019differentiable,chasnov2019convergence,fiez2019convergence} to stabilize the convergence in \emph{general continuous} games, these assumptions are not easily verified/satisfied in MARL, e.g., even in the simplest LQ setting \cite{mazumdar2019policy},   as they depend not only on the model, but also on the policy parameterization. Due to this subtlety, it may be interesting to explore the (global) convergence of policy-based methods for MARL, probably starting with the simple LQ setting, i.e., general-sum LQ games, in analogy to that for the zero-sum counterpart \cite{zhang2019policyb,bu2019global}. 
Such an exploration may also benefit from the recent advances of nonconvex-(non)concave optimization  \cite{lin2018solving,jin2019minmax,nouiehed2019solving}. 

\vspace{6pt}
\noindent{\bf MARL with robustness/safety concerns:} 
\vspace{3pt}
Concerning the challenge of non-unique learning goals in MARL (see \S\ref{subsec:learning_goal}), we believe it is of merit to consider robustness and/or safety constraints in MARL. To the best of our knowledge, this is still a relatively uncharted territory. In fact, safe RL has been recognized as one of the most significant challenges  in the single-agent setting \cite{garcia2015comprehensive}. With more than one agents that may have conflicted objectives, guaranteeing  safety   becomes more involved, as the safety requirement now concerns the coupling of all  agents. One straightforward model is constrained multi-agent MDPs/Markov games, with the constraints characterizing the safety requirement. Learning with provably safety guarantees in this setting is non-trivial, but  necessary for some safety-critical MARL applications as autonomous driving \cite{shalev2016safe} and robotics \cite{kober2013reinforcement}. In addition, it is also natural to think of  robustness against adversarial agents, especially in the decentralized/distributed cooperative MARL settings as in \cite{zhang2018fully,chen2018communication,wai2018multi}, where the adversary may disturb the learning process in an anonymous way -- a common scenario  in distributed systems. Recent development of robust distributed supervised-learning against Byzantine adversaries  \cite{chen2017distributed,yin2018byzantine} may be useful  in this context.  


%\begin{itemize}
%\item Policy gradient for MARL with function approximation, include the simple LQ setting
%\item  Stackelberg equilibrium can be found using Q-learning when the dynamics is controlled only by the leader
%\item MARL in partially observed Stochastic games (POSGs), very hard.
%\item Safe/robust MARL
%\item Model-based MARL with sample complexity analysis include: the early work  \cite{brafman2000near}  proposes polynomial time learning algorithms for single-controller-stochastic games, a special zero-sum MG; \cite{brafman2002r} proposes \text{R-MAX}. R-MAX concerns regret, average-reward settings, 
%\item Neural/Deep MARL theory
%\item preference 
%\end{itemize}